\documentclass{article}
\usepackage[UTF8]{ctex}
\usepackage{geometry}
\usepackage{multirow}
\usepackage{natbib}
\geometry{left=3.18cm,right=3.18cm,top=2.54cm,bottom=2.54cm}
\usepackage{graphicx}
\pagestyle{plain}	
\usepackage{setspace}
\usepackage{enumerate}
\usepackage{caption2}
\usepackage{datetime} %日期
\renewcommand{\today}{\number\year 年 \number\month 月 \number\day 日}
\renewcommand{\captionlabelfont}{\small}
\renewcommand{\captionfont}{\small}
\begin{document}

\begin{figure}
    \centering
    \includegraphics[width=8cm]{upc.png}

    \label{figupc}
\end{figure}

	\begin{center}
		\quad \\
		\quad \\
		\heiti \fontsize{45}{17} \quad \quad \quad 
		\vskip 1.5cm
		\heiti \zihao{2} 《计算科学导论》个人职业规划
	\end{center}
	\vskip 2.0cm
		
	\begin{quotation}
% 	\begin{center}
		\doublespacing
		
        \zihao{4}\par\setlength\parindent{7em}
		\quad 

		学生姓名：\underline{\qquad  汪子霄 \qquad \qquad}

		学\hspace{0.61cm} 号：\underline{\qquad 1907010222\qquad}
		
		专业班级：\underline{\qquad 计科1902 \qquad  }
		
        学\hspace{0.61cm} 院：\underline{计算机科学与技术学院}
% 	\end{center}
		\vskip 1.5cm
		\centering
		\begin{table}[h]
            \centering 
            \zihao{4}
            \begin{tabular}{|c|c|c|c|c|c|c|c|c|}
            % 这里的rl 与表格对应可以看到，姓名是r，右对齐的；学号是l，左对齐的；若想居中，使用c关键字。
                \hline
                \multicolumn{5}{|c|}{分项评价} &\multicolumn{2}{c|}{整体评价}  & 总    分 & 评 阅 教 师\\
                \hline
                自我 & 环境 & 职业 & 实施 & 评估与 & 完整性 & 可行性 &\multirow{2}*{} &\multirow{2}*{}\\
                分析& 分析& 定位 & 方案 & 调整 & 20\% & 20\% & ~&~ \\\            
                10\% & 10\% & 15\% & 15\% & 10\% & &  &~ &~\\
                \cline{1-7} 
                & & & & & & & ~&~ \\
                & & & & & & & ~&~ \\
                \hline      
            \end{tabular}
        \end{table}
		\vskip 2cm
		\today
	\end{quotation}

\thispagestyle{empty}
\newpage
\setcounter{page}{1}
% 在这之前是封面，在这之后是正文
\section{自我分析}
	自我分析，是指对自我理性、深刻、全面的分析，他是自我介绍与自我评价的结合。进行自我分析对每个人来说都是非常有必要的，人在不断的变化、进步，自我的分析也应该不断的更新。古人有曰：“知彼知己，百战不殆”，相比孙老师一定是为了让我们正确知己，向我们颁布这次规划。而作为新晋大学生，我非常感谢孙老师的这次指导，我也认为只有认识了自己，才能对自己的职业做出正确的选择，才能选定适合自己发展的职业生涯路线，才能对自己的职业生涯目标做出最佳抉择。\par

\subsection{自然条件}
我是一名十八岁的男大学生，健康状态保持在非常优秀的水平。我来自山东东营，出自孙武故里广饶的一个普通家庭，拥有着朴素的外貌与不平凡的内心。\par

\subsection{性格分析}
我性格乐观开朗，心里活动丰富；敢于探索有挑战的问题，能够沉着冷静地面对问题，能够在兴趣的驱动下探究问题的本质。我追求上进，能够在帮助他人的过程中发现个人的价值。\par

\subsection{教育与学习经历}
顺利接受九年义务教育，在高中接触信息学竞赛，初次接受到C++语言，虽然只是面向过程的简单问题的解决方法，但它打开了我新知识领域的大门，丰富了高中知识的学习。\par
继高中学习信息学竞赛之后，我在大学选择了计算机科学与技术专业。计算科学导论的学习，让我发现了计算科学的新天地，程序设计的课程，也让我发现了更多富有特点的编程语言（Java，Python）。大学尚未结束，我相信大一构建大体规划后，今后的学习将会更加深入人心。\par

\subsection{工作与社会阅历}
\begin{enumerate}[(1)]
	\item 为学弟学妹讲过C++基础算法课程
	\item 做过初高中衔接班数学助理老师
	\item 乐安公园环保志愿者
	\item 敬老院陪伴老人志愿者
\end{enumerate}
\par

\subsection{知识、技能与经验}
\begin{enumerate}[(1)]
	\item 熟练使用office
	\item 掌握高中的数学物理化学知识
	\item 掌握基本的面向过程的C++代码实现
\end{enumerate}
\par

\subsection{兴趣爱好与特长}
我喜欢听音乐，非常享受在音乐的还要中思考问题与近期的计划。长跑是我锻炼身体的主要方式，我喜欢在奔跑中忘却烦恼。阅读是我放松思想的主要方式，正如孙老师倡导，我时常翻阅《计算科学导论》，在其中发现了我大学的课程大致内容与目前主打方向——网络安全。\par

\section{环境分析}

\subsection{社会环境分析}
习主席深刻指出：“没有网络安全就没有国家安全，就没有经济社会稳定运行，广大人民群众利益也难以得到保障。”作为全球网民最多，联网区域最广，信息化程度与日俱增的网络大国，要确保网络安全，必须构筑好网络安全屏障，确保网络安全责任落到实处。应对网络和信息安全挑战，关键在人。网络空间安全、网络信息安全，在如今大数据运行的今天，我们应培养对自身信息安全基本的防范意识。\par
互联网的发展，种种网络病毒与网络犯罪也随之而来，为了减少和防止该类犯罪给企业和个人带来的隐患，网络安全工程师这一神秘职业逐渐为人们所熟悉。社会对信息安全服务的需求很大，军队、国防、银行、税务、证券、机关、电子商务都急需大批网络安全人才，网络安全工程师已跻身IT新贵之列。在我国，根据国家信息化建设的规模保守估计，2011年国内网络安全专业人才仍存在近百万的巨大缺口，高级的战略人才和专业技术人才尤其匮乏。\par
目前，中国网络安全人才正处在一个起步过程，网络上充斥着大量黑客培训，多属于技巧训练，和真正的网络安全人才培养有本质区别，是完全不同的概念。如果这种局面不进一步改变，网络安全人才的空档只会继续拉大，直接导致国家经济损失加大。\par
作为世界第二大经济体，我国经济仍处于相对高速发展过程中，当今世界正处于百年未有之大变局,经济与生活都在高速发展。\par
如今，5G问题备受政界、商界及学界的普遍关注和广泛讨论,并已成为中美高科技竞争的核心议题，而且5G技术对网络空间安全的影响也是不可忽视的。中国必将需要大量的网络安全工程师，更需要优秀的工程师。\par

\subsection{家庭环境分析}
目前未婚，经济条件一般，基本可以跟随国家的发展，成为小康家庭。我父母都对我充满信心，富有期望，他们希望我可以有自己的方向，有自己的规划，获得自己期望的生活，他们并不会向我过度施压，更多的是时刻鼓励我，为我排忧解难。我的家庭中每个人都秉承艰苦奋斗的精神。\par
\subsection{职业环境分析}
市场需求分析：当今社会更需要的的人才，而人才又分为全才和专才，中国政治稳定，经济持续发展。在全球经济一体化环境中的重要角色。经济发展有强劲的势头，加入WTO后，会有大批的外国企业进入中国市场，中国的企业也将走出国门。\par
职业观察分析：中国的计算机人才不强，而专业人士更少，并且中国科技还不够发达，所以在中国有更多的职业需求，更容易就业。\par
行业分析：计算机软件主流开发技术、软件工程、软件项目过程管理等基本知识与技能、熟练掌握先进的软件开发工具、环境和软件工程管理方法，培养学生系统的软件设计与项目实施能力，胜任软件开发、管理和维护等相关工作的专业性软件
工程高级应用型人才。\par


\subsection{地域与人际环境分析}
我理想的城市坐标南方，而深圳被誉为“中国硅谷”，我也将其作为我的奋斗目标。\par
深圳地处中国华南地区、广东南部、珠江口东岸，东临大亚湾和大鹏湾，西濒珠江口和伶仃洋，南隔深圳河与香港相连，所处纬度较低，属亚热带海洋性气候。由于深受季风的影响，夏季盛行偏东南风，时有季风低压、热带气旋光顾，高温多雨；其余季节盛行东北季风，天气较为干燥，气候温和。他是粤港澳大湾区四大中心城市之一 、国际性综合交通枢纽、国际科技产业创新中心 、中国三大全国性金融中心之一 ，并全力建设全球海洋中心城市、中国特色社会主义先行示范区、综合性国家科学中心。深圳水陆空铁口岸俱全，是中国拥有口岸数量最多、出入境人员最多、车流量最大的口岸城市。\par
深圳是中国经济中心城市之一，经济总量长期位列中国大陆城市第四位，是中国大陆经济效益最好的城市之一。英国《经济学人》2012年“全球最具经济竞争力城市”榜单上，深圳位居第二。中国优秀企业华为更是坐标深圳，城市的优秀决定人才的优秀，我相信与优秀的人竞争可以得到优秀的人际关系与提升程度。\par 

\section{职业定位}

\subsection{行业领域定位与理由}
从事网络信息安全领域。\par
网络工程是国家战略工程，网络安全问题关系到国家的安全与社会的稳定，在网络信息技术高速发展的今天，在全球化进程的不断加速中，网络安全的重要性被日益放大，解决存在的安全问题变得非常迫切。其重要性正随着全球信息化步伐的加快而变得越来越重要，安全问题需要迫切解决。从微观角度看，网络安全关系到一个单位的发展。在企事业单位中，无论是科研数据还是财务数据，甚至是人力资源数据，都需要通过计算机网络来处理，特别是自动化办公的推广，企事业单位的很多工作都已经离不开网络，但是网络又离不开安全支撑，一旦出现安全问题，其后果是非常严重的。在网络工程中，安全技术的运用占据着非常重要的地位。\par
目前，许多企事业单位的业务依赖于信息系统安全运行，信息安全重要性日益凸显。信息已经成为各企事业单位中重要资源，也是一种重要的“无形财富”，在未来竞争中谁获取信息优势，谁就掌握了竞争的主动权。信息安全已成为影响国家安全、经济发展、社会稳定、个人利害的重大关键问题。\par

\subsection{职业岗位起点定位与理由}
网络安全工程师\par
对网络安全的探索是我的兴趣驱动，可以让我的效率事半功倍。我认为以网络安全为起点，可以进一步夯实网络基础，为未来方向提供一定的实力基础。\par

\subsection{职业目标与可行性分析}
\begin{enumerate}[(1)]
	\item 短期目标（大学4年）
	\begin{table}[h]
		\centering 
		\zihao{5}
		\begin{tabular}{|c|c|c|}
			% 这里的rl 与表格对应可以看到，姓名是r，右对齐的；学号是l，左对齐的；若想居中，使用c关键字。
			\hline
			年级 & 主要方向 & 课外计划\\
			\hline
			大学 & 高等数学、离散数学、 & 学习Java、HTML、JavaScrip\\
			一年级 &  线性代数、C++、python & 巩固ACM知识，准备省赛 \\ 
			\hline
			大学 & 掌握专业课知识 & 考虑大学生创新与数学建模竞赛\\
			二年级 & 备考CET-4 &  加强ACM代码能力，争取区域赛\\
			\hline
			大学 & 学习信息安全等专业知识 & 学习JavaServer Pages、Linux \\
			三年级 & 备考CET-6 & 在继续ACM与其他项目中做出选择 \\
			\hline
			大学 & 在备考研究生与 & 研究大学实力或者 \\
			四年级 & 工作之间做出选择 & 了解企业的招聘信息 \\
			\hline
		\end{tabular}
	\end{table}
	\item 中长期目标（5-10年）\par
	保持自主学习的能力。\par
	研究生方向选择网络安全，发表自己对网络安全理解的论文 \par
	就业选择网络安全工程师，不断加强自己的能力，不局限与单一性学习。\par
\end{enumerate}


\section{实施方案}
\begin{enumerate}[1、]
	\item 充分利用乐于探索的精神，保持并发扬对数学的热爱。
	\item 克服拖延症，提升英语水平，利用琐碎时间阅读英文原著。
	\item 通过询问老师或网上资源的途径自主训练，目标能够担负起小型网络信息安全工作。对网络信息安全有较为完整的认识，掌握电脑安全防护、网站安全、电子邮件安全、Internet网络安全部署、操作系统安全配置、恶意代码防护、常用软件安全设置、防火墙的应用等技能。
	\item 在深入专业训练专业训练，能完善和优化企业信息安全制度和流程。信息安全工作符合特定的规范要求，能够对系统中安全措施的实施进行了跟踪和验证，能够建立起立体式、纵深的安全防护系统，部署安全监控机制，对未知的安全威胁能够进行预警和追踪。
	\item 利用音乐与跑步发泄情绪，善于发泄学会发泄
\end{enumerate}
\par 

\section{评估与调整}

\subsection{评估时间}
每学期评估一次。\par 

\subsection{评估内容}
计算机语言的掌握水平，竞赛成果，英语水平，课程学分绩。\par

\subsection{调整原则}
自身情况的匹配性、与环境的适应性、操作实施的可行性等。\par


\end{document}
